% !Mode:: "TeX:UTF-8"
% !TEX root = ../../main.tex
% ============================================================
% TOPIC 2. DETERMINANTS OF ORDER TWO AND THREE
% ============================================================
\mashgulot{2}{The concept of a determinant. Determinants of the second and third order}
\begin{tarif}
With every square matrix of order $n$ we associate a single number called its
\textbf{determinant}. The determinant of a matrix $A$ is denoted by $\det A$, $|A|$
or $\Delta$:
\[\kf{\det A = |A| =
\begin{vmatrix}
a_{11} & a_{12} & \cdots & a_{1n} \\
a_{21} & a_{22} & \cdots & a_{2n} \\
\vdots & \vdots & \ddots & \vdots \\
a_{n1} & a_{n2} & \cdots & a_{nn}
\end{vmatrix}.}\]
Determinants are defined \textbf{only for square matrices}.
\end{tarif}
\subsection{Second-order determinants}
\vspace{3pt}
\begin{tarif}
The determinant of a $2\times 2$ matrix is given by
\[\kf{\det A =
\begin{vmatrix}
a_{11} & a_{12} \\
a_{21} & a_{22}
\end{vmatrix}
= a_{11}\cdot a_{22} - a_{12}\cdot a_{21}.}\]
In words, we subtract the product of the entries on the \textbf{secondary diagonal}
from the product of the entries on the \textbf{main diagonal}.
\end{tarif}
\begin{misol}
$\begin{vmatrix} 3 & -2 \\ 5 & 4 \end{vmatrix}
= 3\cdot 4 - (-2)\cdot 5 = 12 + 10 = 22.$
\end{misol}
\begin{misol}
$\begin{vmatrix} -1 & 7 \\ 2 & -3 \end{vmatrix}
= (-1)\cdot(-3) - 7\cdot 2 = 3 - 14 = -11.$
\end{misol}
\begin{misol}
$\begin{vmatrix} 0 & 5 \\ -4 & 1 \end{vmatrix}
= 0\cdot 1 - 5\cdot(-4) = 0 + 20 = 20.$
\end{misol}
\subsection{Third-order determinants}
\vspace{3pt}
\begin{tarif}
Two methods for evaluating a $3\times 3$ determinant are in common use:
\textbf{Sarrus’ rule} and \textbf{expansion along the first row} (the Laplace expansion).
\end{tarif}
\textbf{Sarrus’ rule (the rule of diagonals).}
The scheme behind Sarrus’ rule for a $3\times3$ determinant is shown in Figure~1.
\begin{center}
\includegraphics{sarrus-sxema.pdf}
\end{center}
\begin{center}
    Figure 1.
\end{center}

In practice one copies the first two columns to the right of the matrix. In the
resulting $3\times 5$ array the products along the three diagonals running down
from left to right are taken with the sign “$+$” and \textbf{added}, while the products
along the three diagonals running up from left to right are taken with the sign
“$-$” and \textbf{subtracted}. The same products can be read off the triangle scheme
of Figure~2.
\begin{figure}[htbp]
  \centering
  \includegraphics[width=\linewidth]{uchburchaklar-sxemasi.pdf}
\end{figure}
\begin{center}
    Figure 2.
\end{center}
Either scheme leads to the formula
\[\kf{\begin{vmatrix}
a_{11} & a_{12} & a_{13} \\
a_{21} & a_{22} & a_{23} \\
a_{31} & a_{32} & a_{33}
\end{vmatrix}
= a_{11}\times a_{22}\times a_{33} + a_{12}\times a_{23}\times a_{31} + a_{13}\times a_{21}\times a_{32}-}\]
\[\kf{{} - a_{13}\times a_{22}\times a_{31} - a_{12}\times a_{21}\times a_{33} - a_{11}\times a_{23}\times a_{32}.}\]
\vspace{3pt}
\begin{misol}
We evaluate by Sarrus’ rule
\[\kf{\begin{vmatrix}
2 & 1 & 3 \\
0 & -1 & 2 \\
4 & 3 & 1
\end{vmatrix}.}\]
The diagonals with the sign “$+$” give
$2\cdot(-1)\cdot 1 + 1\cdot 2\cdot 4 + 3\cdot 0\cdot 3 = -2 + 8 + 0 = 6$,
and those with the sign “$-$” give
$3\cdot(-1)\cdot 4 + 1\cdot 0\cdot 1 + 2\cdot 2\cdot 3 = -12 + 0 + 12 = 0$.
Hence
\[\kf{\det A = 6 - 0 = 6.}\]
\end{misol}
\subsection{Basic properties of determinants}
The following properties hold for determinants of any order.
\begin{enumerate}[label=\arabic*)]
\item \textbf{Transposition:} $\det A^{T} = \det A$; interchanging the roles of rows
and columns does not change the value of a determinant.
\item \textbf{Change of sign:} interchanging two rows (or two columns) changes the
sign of the determinant.
\item \textbf{Zero row (column):} if all entries of some row (or column) are zero,
the determinant equals zero.
\item \textbf{Common factor:} if all entries of one row (or column) are multiplied
by a number $k$, the determinant is multiplied by $k$. In other words, a common
factor of a row can be taken outside the determinant sign. Since $kA$ has every one
of its $n$ rows multiplied by $k$,
\[\kf{\det(kA) = k^{n}\det A \quad (\text{for a matrix of order } n).}\]
\item \textbf{Equal rows:} if two rows (or columns) of a determinant are equal,
the determinant equals zero.
\item \textbf{Proportional rows:} if one row is proportional to another, the
determinant equals zero.
\item \textbf{Adding a multiple of a row:} adding to one row a multiple $k$ of
another row does not change the value of the determinant.
\item \textbf{Determinant of a product:}
$\det(AB)=\det A \cdot \det B$.
\end{enumerate}
\vspace{3pt}
\begin{misol}
\textbf{(Property 4).}
\[\kf{\begin{vmatrix}
6 & 3 \\
4 & 2
\end{vmatrix}
= 3\cdot\begin{vmatrix}2 & 1 \\ 4 & 2\end{vmatrix}
= 3\cdot(2\cdot 2 - 1\cdot 4) = 3\cdot 0 = 0.}\]
Here we took the common factor 3 out of the first row. In the new determinant the
second row is twice the first (Property 6), so the determinant is zero.
\end{misol}
\begin{misol}
\textbf{(Property 7).} We use the properties to evaluate
\[\kf{\begin{vmatrix}
1 & 2 & 3 \\
4 & 5 & 6 \\
7 & 8 & 9
\end{vmatrix}
\stackrel{S_2 - 4S_1}{=}
\begin{vmatrix}
1 & 2 & 3 \\
0 & -3 & -6 \\
7 & 8 & 9
\end{vmatrix}
\stackrel{S_3 - 7S_1}{=}
\begin{vmatrix}
1 & 2 & 3 \\
0 & -3 & -6 \\
0 & -6 & -12
\end{vmatrix}.}\]
Here $S_i$ denotes the $i$th row. The second and third rows are now proportional
($S_3 = 2S_2$), and therefore $\det A = 0$.
\end{misol}
\subsection{Determinants in information security}
Determinants appear naturally in information security and cryptography.

\textbf{Recovering a key from a linear system.} If the system of equations written
for the unknown entries of a key has a nonzero determinant, the key is determined
uniquely. This is exactly what makes linear ciphers vulnerable to known-plaintext
attacks.

\textbf{The Hill cipher.} A key matrix $K$ is admissible only if its determinant
satisfies $\gcd(\det K, 26) = 1$. When this condition fails, the inverse matrix
$K^{-1} \pmod{26}$ does not exist and decryption is impossible. Two cases are
typical. If $\det K \equiv 0 \pmod{26}$, the matrix is \textbf{singular} modulo 26 and
cannot serve as a key at all. If $\gcd(\det K, 26) \neq 1$, for instance when
$\det K$ is even or divisible by 13, the number $\det K$ has no inverse modulo 26.
Choosing a key therefore requires evaluating determinants quickly, and for small
keys this rests directly on the formulas for determinants of order two and three.
\vspace{3pt}
\begin{misol}
For $K=\begin{pmatrix}3 & 5\\ 2 & 7\end{pmatrix}$ we have
$\det K = 3\cdot 7 - 5\cdot 2 = 21-10=11$ and $\gcd(11,26)=1$, so $K$ is an admissible
Hill cipher key. For $K'=\begin{pmatrix}4 & 6\\ 2 & 3\end{pmatrix}$ we get
$\det K' = 12-12=0$, and this matrix is \textbf{not admissible} as a key.
\end{misol}
\begin{misol}
Compute $\det K$ for $K=\begin{pmatrix}5 & 3\\ 2 & 7\end{pmatrix}$ and decide whether
$K$ can be used as a Hill cipher key.

\textbf{Solution.} $\det K = 5\cdot 7 - 3\cdot 2 = 35 - 6 = 29$.
By the Euclidean algorithm, $\gcd(29, 26) = \gcd(26, 3) = \gcd(3, 2) = \gcd(2, 1) = 1$.
The condition $\gcd(\det K, 26) = 1$ holds, so $K$ is an admissible Hill cipher key.
\end{misol}
\textbf{Checking that an inverse exists.}
Most cryptographic systems need to “undo” a key: to decrypt a message, to verify a
signature, and so on. This requires an inverse matrix, which exists only when
$\det A \neq 0$. Computing the determinant is therefore the first check in the
design of such a system.

\textbf{Determinants in error-correcting codes.}
In coding theory the rows of the generator and parity-check matrices must be
linearly independent. To verify this one computes the determinants of suitable
square submatrices (minors). If the required minors are nonzero, the code is
correctly constructed and is able to detect errors.

\textbf{Adjacency matrices in network security.}
In the analysis of network topology, the determinant and other spectral
characteristics (eigenvalues) of the adjacency matrix are used to assess how
well connected and how resilient a network is. For example, by
\textbf{Kirchhoff’s matrix-tree theorem} the number of spanning trees of a network
equals any cofactor of its Laplacian matrix, that is, a certain determinant. This
number measures how many backup routes the network has.

\subsection{Computing with Python, Maple and Mathcad}
Determinants of numerical matrices are easily checked by computer. The examples
below evaluate the third-order determinant from the Sarrus example and test the
Hill cipher key $K=\begin{pmatrix}5&3\\2&7\end{pmatrix}$.
In Python we use SymPy, which works with exact integers:
\begin{Verbatim}[fontsize=\footnotesize,frame=single,label={\small Python}]
from math import gcd
from sympy import Matrix
A = Matrix([[2, 1, 3], [0, -1, 2], [4, 3, 1]])
print(A.det())
K = Matrix([[5, 3], [2, 7]])
d = K.det()
print(d, gcd(int(d), 26))
\end{Verbatim}
\begin{Verbatim}[fontsize=\footnotesize,frame=single,label={\small Output}]
6
29 1
\end{Verbatim}
In Maple the same computation uses the \texttt{LinearAlgebra} package:
\begin{Verbatim}[fontsize=\footnotesize,frame=single,label={\small Maple}]
with(LinearAlgebra):
A := Matrix([[2, 1, 3], [0, -1, 2], [4, 3, 1]]):
Determinant(A);
                               6
K := Matrix([[5, 3], [2, 7]]):
igcd(Determinant(K), 26);
                               1
\end{Verbatim}
In Mathcad Prime a matrix is entered through the matrix template (Ctrl+M), and
the determinant is obtained with the operator $|A|$ from the Matrices/Vectors
toolbar (the function \texttt{gcd} is built in):
\begin{Verbatim}[fontsize=\footnotesize,frame=single,label={\small Mathcad}]
     | 2   1  3 |
A := | 0  -1  2 |        |A| = 6
     | 4   3  1 |

K := | 5  3 |            |K| = 29
     | 2  7 |
gcd(|K|, 26) = 1
\end{Verbatim}

\subsection{Exercises on determinants}
\bigskip
\emph{1. Evaluate the second-order determinants.}
\bigskip
\begin{multicols}{3}\raggedright
\begin{enumerate}[label=\arabic*.,  leftmargin=*]
\item $\begin{vmatrix} 3 & -2 \\ 5 & 4 \end{vmatrix}$
\item $\begin{vmatrix} -1 & 7 \\ 2 & -3 \end{vmatrix}$
\item $\begin{vmatrix} 0 & 5 \\ -4 & 1 \end{vmatrix}$
\item $\begin{vmatrix} 6 & -3 \\ -2 & 8 \end{vmatrix}$
\item $\begin{vmatrix} -5 & 1 \\ 3 & -7 \end{vmatrix}$
\item $\begin{vmatrix} 2 & 4 \\ -1 & 6 \end{vmatrix}$
\item $\begin{vmatrix} 7 & 0 \\ -3 & 5 \end{vmatrix}$
\item $\begin{vmatrix} -4 & 2 \\ 9 & -1 \end{vmatrix}$
\item $\begin{vmatrix} 1 & -8 \\ 3 & 2 \end{vmatrix}$
\item $\begin{vmatrix} -6 & 3 \\ 4 & -5 \end{vmatrix}$
\item $\begin{vmatrix} 8 & -1 \\ 2 & 7 \end{vmatrix}$
\item $\begin{vmatrix} -3 & 9 \\ 1 & -4 \end{vmatrix}$
\item $\begin{vmatrix} 5 & 2 \\ -6 & 3 \end{vmatrix}$
\item $\begin{vmatrix} -2 & -5 \\ 4 & 1 \end{vmatrix}$
\item $\begin{vmatrix} 4 & 7 \\ -3 & -2 \end{vmatrix}$
\item $\begin{vmatrix} 1 & -1 \\ -1 & 1 \end{vmatrix}$
\item $\begin{vmatrix} -7 & 4 \\ 5 & 3 \end{vmatrix}$
\item $\begin{vmatrix} 9 & -6 \\ 2 & 8 \end{vmatrix}$
\item $\begin{vmatrix} 3 & 5 \\ -8 & -4 \end{vmatrix}$
\item $\begin{vmatrix} -4 & -3 \\ 7 & 6 \end{vmatrix}$
\item $\begin{vmatrix} 2 & -9 \\ 1 & 5 \end{vmatrix}$
\item $\begin{vmatrix} 6 & 1 \\ -5 & -8 \end{vmatrix}$
\item $\begin{vmatrix} -3 & 4 \\ 8 & -2 \end{vmatrix}$
\item $\begin{vmatrix} 10 & -7 \\ 3 & 1 \end{vmatrix}$
\item $\begin{vmatrix} -8 & 5 \\ 6 & -9 \end{vmatrix}$
\item $\begin{vmatrix} 5 & -4 \\ -2 & 3 \end{vmatrix}$
\item $\begin{vmatrix} -1 & 6 \\ 7 & -5 \end{vmatrix}$
\item $\begin{vmatrix} 4 & 3 \\ -9 & 2 \end{vmatrix}$
\item $\begin{vmatrix} -6 & -2 \\ 1 & 8 \end{vmatrix}$
\item $\begin{vmatrix} 7 & -5 \\ -4 & 3 \end{vmatrix}$
\end{enumerate}
\end{multicols}
\emph{2. Evaluate the third-order determinants by Sarrus’ rule.}
\bigskip
\begin{multicols}{3}\raggedright
\begin{enumerate}[label=\arabic*.,  leftmargin=*]
\item $\begin{vmatrix} 2 & 1 & 3 \\ 0 & -1 & 2 \\ 4 & 3 & 1 \end{vmatrix}$
\item $\begin{vmatrix} 1 & -3 & 2 \\ 5 & 0 & -1 \\ -2 & 4 & 3 \end{vmatrix}$
\item $\begin{vmatrix} 3 & 2 & -1 \\ 1 & 4 & 0 \\ -2 & 5 & 3 \end{vmatrix}$
\item $\begin{vmatrix} -1 & 2 & 3 \\ 4 & 0 & -2 \\ 1 & 5 & 1 \end{vmatrix}$
\item $\begin{vmatrix} 5 & -1 & 2 \\ 3 & 4 & -3 \\ 0 & 2 & 1 \end{vmatrix}$
\item $\begin{vmatrix} 2 & 0 & -4 \\ -1 & 3 & 5 \\ 6 & -2 & 1 \end{vmatrix}$
\item $\begin{vmatrix} -3 & 1 & 2 \\ 0 & 5 & -4 \\ 4 & -1 & 3 \end{vmatrix}$
\item $\begin{vmatrix} 1 & 2 & -3 \\ -4 & 0 & 1 \\ 3 & 5 & 2 \end{vmatrix}$
\item $\begin{vmatrix} 4 & -2 & 1 \\ 2 & 3 & -5 \\ -1 & 0 & 6 \end{vmatrix}$
\item $\begin{vmatrix} -2 & 5 & 0 \\ 1 & -3 & 4 \\ 3 & 2 & -1 \end{vmatrix}$
\item $\begin{vmatrix} 0 & 3 & -2 \\ 5 & 1 & 4 \\ -3 & -4 & 2 \end{vmatrix}$
\item $\begin{vmatrix} 3 & -4 & 1 \\ -2 & 0 & 5 \\ 1 & 3 & -3 \end{vmatrix}$
\item $\begin{vmatrix} 1 & 0 & 5 \\ -3 & 2 & -1 \\ 4 & -4 & 3 \end{vmatrix}$
\item $\begin{vmatrix} -5 & 2 & 1 \\ 3 & -1 & 0 \\ 2 & 4 & -3 \end{vmatrix}$
\item $\begin{vmatrix} 2 & -1 & 4 \\ 0 & 3 & -2 \\ 5 & 1 & -1 \end{vmatrix}$
\item $\begin{vmatrix} 4 & 1 & -3 \\ -1 & -2 & 5 \\ 3 & 0 & 2 \end{vmatrix}$
\item $\begin{vmatrix} -2 & 3 & 1 \\ 4 & -5 & 2 \\ 0 & 1 & -4 \end{vmatrix}$
\item $\begin{vmatrix} 1 & -2 & 0 \\ 3 & 4 & -1 \\ -5 & 2 & 3 \end{vmatrix}$
\item $\begin{vmatrix} 6 & 0 & -3 \\ -1 & 2 & 5 \\ 4 & -3 & 1 \end{vmatrix}$
\item $\begin{vmatrix} -3 & 4 & 2 \\ 1 & -1 & -5 \\ 5 & 3 & 0 \end{vmatrix}$
\item $\begin{vmatrix} 0 & -5 & 3 \\ 2 & 1 & -4 \\ -1 & 6 & 2 \end{vmatrix}$
\item $\begin{vmatrix} 5 & 3 & -2 \\ -4 & 0 & 1 \\ 2 & -1 & 4 \end{vmatrix}$
\item $\begin{vmatrix} -1 & 4 & -3 \\ 0 & 2 & 5 \\ 3 & -1 & 1 \end{vmatrix}$
\item $\begin{vmatrix} 2 & -3 & 5 \\ 4 & 1 & 0 \\ -1 & -2 & 3 \end{vmatrix}$
\item $\begin{vmatrix} 3 & 0 & 1 \\ -2 & 5 & -4 \\ 4 & 1 & 2 \end{vmatrix}$
\item $\begin{vmatrix} -4 & 1 & 3 \\ 2 & -3 & 0 \\ 5 & 4 & -2 \end{vmatrix}$
\item $\begin{vmatrix} 1 & 5 & -1 \\ -3 & 0 & 4 \\ 2 & -2 & 3 \end{vmatrix}$
\item $\begin{vmatrix} 0 & -1 & 4 \\ 3 & 2 & -5 \\ -4 & 1 & 0 \end{vmatrix}$
\item $\begin{vmatrix} -2 & 0 & 3 \\ 1 & -4 & 5 \\ 4 & 2 & -1 \end{vmatrix}$
\item $\begin{vmatrix} 3 & -2 & 0 \\ -1 & 5 & 4 \\ 2 & 0 & -3 \end{vmatrix}$
\end{enumerate}
\end{multicols}

\emph{3. Evaluate the third-order determinants by expansion along the first row.}
\bigskip
\begin{multicols}{3}\raggedright
\begin{enumerate}[label=\arabic*.,  leftmargin=*]
\item $\begin{vmatrix} 1 & 3 & -2 \\ 4 & 0 & 5 \\ -1 & 2 & 3 \end{vmatrix}$
\item $\begin{vmatrix} -2 & 1 & 4 \\ 3 & -5 & 0 \\ 1 & 2 & -3 \end{vmatrix}$
\item $\begin{vmatrix} 5 & -3 & 1 \\ 0 & 2 & -4 \\ -1 & 3 & 5 \end{vmatrix}$
\item $\begin{vmatrix} 3 & 0 & -1 \\ -2 & 4 & 3 \\ 5 & -1 & 2 \end{vmatrix}$
\item $\begin{vmatrix} -4 & 2 & 0 \\ 1 & -3 & 5 \\ 3 & 4 & -2 \end{vmatrix}$
\item $\begin{vmatrix} 0 & 5 & -3 \\ -2 & 1 & 4 \\ 4 & -1 & 2 \end{vmatrix}$
\item $\begin{vmatrix} 2 & -4 & 3 \\ 5 & 0 & -1 \\ -3 & 2 & 4 \end{vmatrix}$
\item $\begin{vmatrix} -1 & 3 & 5 \\ 4 & -2 & 0 \\ 0 & 1 & -3 \end{vmatrix}$
\item $\begin{vmatrix} 6 & 1 & -2 \\ -3 & 0 & 4 \\ 2 & 5 & 1 \end{vmatrix}$
\item $\begin{vmatrix} 1 & -1 & 0 \\ 0 & 3 & -5 \\ -2 & 4 & 1 \end{vmatrix}$
\item $\begin{vmatrix} -3 & 0 & 4 \\ 2 & 5 & -1 \\ 1 & -2 & 3 \end{vmatrix}$
\item $\begin{vmatrix} 4 & -2 & 5 \\ -1 & 3 & 0 \\ 2 & 0 & -4 \end{vmatrix}$
\item $\begin{vmatrix} 0 & 4 & 1 \\ 5 & -3 & 2 \\ -2 & 1 & -5 \end{vmatrix}$
\item $\begin{vmatrix} -5 & 1 & 0 \\ 3 & -2 & 4 \\ 0 & 5 & -1 \end{vmatrix}$
\item $\begin{vmatrix} 2 & 0 & -3 \\ -4 & 1 & 5 \\ 3 & -2 & 0 \end{vmatrix}$
\item $\begin{vmatrix} 1 & -5 & 2 \\ 0 & 4 & -3 \\ -3 & 1 & 0 \end{vmatrix}$
\item $\begin{vmatrix} -2 & 4 & -1 \\ 5 & 0 & 3 \\ 1 & -3 & 2 \end{vmatrix}$
\item $\begin{vmatrix} 3 & -1 & 4 \\ 0 & 2 & -5 \\ -4 & 3 & 1 \end{vmatrix}$
\item $\begin{vmatrix} -1 & 0 & 3 \\ 2 & -4 & 1 \\ 5 & 3 & -2 \end{vmatrix}$
\item $\begin{vmatrix} 4 & 3 & 0 \\ -2 & -1 & 5 \\ 0 & 4 & -3 \end{vmatrix}$
\item $\begin{vmatrix} 0 & -2 & 5 \\ 3 & 1 & -4 \\ -1 & 0 & 2 \end{vmatrix}$
\item $\begin{vmatrix} 5 & 1 & -4 \\ -3 & 0 & 2 \\ 4 & -5 & 1 \end{vmatrix}$
\item $\begin{vmatrix} -4 & 3 & 0 \\ 1 & -1 & -2 \\ 0 & 5 & 4 \end{vmatrix}$
\item $\begin{vmatrix} 2 & 0 & 1 \\ -5 & 3 & -4 \\ 1 & -1 & 5 \end{vmatrix}$
\item $\begin{vmatrix} -1 & -3 & 2 \\ 4 & 0 & -5 \\ 3 & 2 & 1 \end{vmatrix}$
\item $\begin{vmatrix} 3 & 2 & -5 \\ 0 & -4 & 1 \\ -2 & 3 & 0 \end{vmatrix}$
\item $\begin{vmatrix} -3 & 5 & 1 \\ 2 & 0 & -4 \\ 0 & -1 & 3 \end{vmatrix}$
\item $\begin{vmatrix} 1 & -4 & 0 \\ 0 & 5 & -2 \\ 3 & -1 & 4 \end{vmatrix}$
\item $\begin{vmatrix} -5 & 0 & 2 \\ 4 & -3 & 1 \\ -1 & 2 & 0 \end{vmatrix}$
\item $\begin{vmatrix} 0 & 1 & -4 \\ -2 & 3 & 5 \\ 4 & 0 & -1 \end{vmatrix}$
\end{enumerate}
\end{multicols}

\emph{4. Use the properties of determinants to evaluate the following determinants.}
\bigskip
\begin{multicols}{3}\raggedright
\begin{enumerate}[label=\arabic*.,  leftmargin=*]
\item $\begin{vmatrix} 2 & 4 \\ 3 & 6 \end{vmatrix}$
\item $\begin{vmatrix} -5 & 10 \\ 1 & -2 \end{vmatrix}$
\item $\begin{vmatrix} 3 & 0 & -6 \\ 1 & 2 & 4 \\ -2 & 1 & 5 \end{vmatrix}$
\item $\begin{vmatrix} 1 & 2 & 3 \\ 4 & 5 & 6 \\ 7 & 8 & 9 \end{vmatrix}$
\item $\begin{vmatrix} 5 & -5 \\ -3 & 3 \end{vmatrix}$
\item $\begin{vmatrix} 2 & 1 & -3 \\ 0 & 0 & 0 \\ 4 & -2 & 5 \end{vmatrix}$
\item $\begin{vmatrix} 1 & 3 & 5 \\ 2 & 6 & 10 \\ 4 & -1 & 3 \end{vmatrix}$
\item $\begin{vmatrix} 6 & -3 & 9 \\ 2 & 4 & -1 \\ 0 & 5 & 3 \end{vmatrix}$
\item $\begin{vmatrix} -4 & 8 \\ 7 & -14 \end{vmatrix}$
\item $\begin{vmatrix} 3 & 1 & 2 \\ 6 & 2 & 4 \\ 1 & -3 & 5 \end{vmatrix}$
\item $\begin{vmatrix} 1 & 0 & 0 \\ 0 & 5 & 0 \\ 0 & 0 & -3 \end{vmatrix}$
\item $\begin{vmatrix} 2 & -1 & 3 \\ 2 & -1 & 3 \\ 4 & 0 & -2 \end{vmatrix}$
\item $\begin{vmatrix} -2 & 6 \\ 1 & -3 \end{vmatrix}$
\item $\begin{vmatrix} 0 & 0 & 0 \\ 3 & -1 & 5 \\ 2 & 4 & -7 \end{vmatrix}$
\item $\begin{vmatrix} 1 & 2 & 3 \\ 2 & 4 & 6 \\ -1 & 3 & 5 \end{vmatrix}$
\item $\begin{vmatrix} 10 & 5 \\ -4 & -2 \end{vmatrix}$
\item $\begin{vmatrix} 3 & -1 & 2 \\ 0 & 4 & -5 \\ 6 & -2 & 4 \end{vmatrix}$
\item $\begin{vmatrix} 1 & 1 & 1 \\ 2 & 2 & 2 \\ 3 & 4 & 5 \end{vmatrix}$
\item $\begin{vmatrix} 7 & 0 \\ 0 & -3 \end{vmatrix}$
\item $\begin{vmatrix} 4 & 2 & -6 \\ 0 & 3 & 1 \\ -2 & -1 & 3 \end{vmatrix}$
\item $\begin{vmatrix} 5 & -10 & 15 \\ 1 & 2 & -3 \\ 0 & 4 & 1 \end{vmatrix}$
\item $\begin{vmatrix} -3 & 0 & 0 \\ 1 & 2 & 0 \\ -4 & 5 & 7 \end{vmatrix}$
\item $\begin{vmatrix} 1 & -3 & 5 \\ -2 & 6 & -10 \\ 3 & 1 & 4 \end{vmatrix}$
\item $\begin{vmatrix} 8 & 4 \\ -6 & -3 \end{vmatrix}$
\item $\begin{vmatrix} 2 & 0 & 1 \\ 0 & -1 & 0 \\ 3 & 0 & 4 \end{vmatrix}$
\item $\begin{vmatrix} 1 & 2 & 3 \\ 3 & 6 & 9 \\ 2 & -1 & 4 \end{vmatrix}$
\item $\begin{vmatrix} -1 & 5 \\ 2 & -10 \end{vmatrix}$
\item $\begin{vmatrix} 2 & 3 & -1 \\ 4 & 6 & -2 \\ 1 & 5 & 3 \end{vmatrix}$
\item $\begin{vmatrix} 0 & 3 & 0 \\ -1 & 0 & 2 \\ 0 & 5 & 0 \end{vmatrix}$
\item $\begin{vmatrix} 12 & -8 & 4 \\ 3 & 1 & -2 \\ 0 & 5 & 6 \end{vmatrix}$
\end{enumerate}
\end{multicols}
\emph{5. Compute $\det K$, check the condition $\gcd(\det K, 26) = 1$, and decide whether the matrix $K$ can be used as a Hill cipher key.}
\bigskip
\begin{multicols}{3}\raggedright
\begin{enumerate}[label=\arabic*.,  leftmargin=*]
\item $K = \begin{pmatrix} 3 & 5 \\ 2 & 7 \end{pmatrix}$
\item $K = \begin{pmatrix} 4 & 6 \\ 2 & 3 \end{pmatrix}$
\item $K = \begin{pmatrix} 5 & 8 \\ 3 & 5 \end{pmatrix}$
\item $K = \begin{pmatrix} 7 & 2 \\ 1 & 3 \end{pmatrix}$
\item $K = \begin{pmatrix} 9 & 4 \\ 5 & 7 \end{pmatrix}$
\item $K = \begin{pmatrix} 2 & 3 \\ 4 & 6 \end{pmatrix}$
\item $K = \begin{pmatrix} 1 & 2 \\ 3 & 5 \end{pmatrix}$
\item $K = \begin{pmatrix} 6 & 1 \\ 5 & 4 \end{pmatrix}$
\item $K = \begin{pmatrix} 8 & 3 \\ 7 & 9 \end{pmatrix}$
\item $K = \begin{pmatrix} 10 & 13 \\ 3 & 4 \end{pmatrix}$
\item $K = \begin{pmatrix} 11 & 5 \\ 2 & 9 \end{pmatrix}$
\item $K = \begin{pmatrix} 4 & 9 \\ 6 & 13 \end{pmatrix}$
\item $K = \begin{pmatrix} 3 & 3 \\ 2 & 5 \end{pmatrix}$
\item $K = \begin{pmatrix} 7 & 8 \\ 11 & 11 \end{pmatrix}$
\item $K = \begin{pmatrix} 1 & 1 \\ 0 & 1 \end{pmatrix}$
\item $K = \begin{pmatrix} 15 & 7 \\ 4 & 2 \end{pmatrix}$
\item $K = \begin{pmatrix} 5 & 3 \\ 7 & 9 \end{pmatrix}$
\item $K = \begin{pmatrix} 12 & 5 \\ 1 & 3 \end{pmatrix}$
\item $K = \begin{pmatrix} 6 & 4 \\ 3 & 2 \end{pmatrix}$
\item $K = \begin{pmatrix} 25 & 1 \\ 3 & 7 \end{pmatrix}$
\item $K = \begin{pmatrix} 2 & 11 \\ 1 & 6 \end{pmatrix}$
\item $K = \begin{pmatrix} 13 & 0 \\ 5 & 2 \end{pmatrix}$
\item $K = \begin{pmatrix} 9 & 7 \\ 4 & 5 \end{pmatrix}$
\item $K = \begin{pmatrix} 8 & 6 \\ 3 & 5 \end{pmatrix}$
\item $K = \begin{pmatrix} 3 & 2 \\ 1 & 1 \end{pmatrix}$
\item $K = \begin{pmatrix} 10 & 9 \\ 7 & 6 \end{pmatrix}$
\item $K = \begin{pmatrix} 4 & 5 \\ 3 & 4 \end{pmatrix}$
\item $K = \begin{pmatrix} 14 & 3 \\ 5 & 1 \end{pmatrix}$
\item $K = \begin{pmatrix} 2 & 7 \\ 5 & 3 \end{pmatrix}$
\item $K = \begin{pmatrix} 11 & 8 \\ 3 & 7 \end{pmatrix}$
\end{enumerate}
\end{multicols}
